% !PRINTED 509
\weberpage{509}

%<<CONT
so kommen nur \gesperrt{eigentliche} Lösungen von (8) in Betracht
(\S\,114). Jede eigentliche Lösung führt aber nach (9) zu einem
Wertsystem $\alpha, \beta$, und durch die Kongruenzen

\weq{10}{c\alpha \equiv Cx, \qquad c\beta \equiv \dfrac{Bx-y}{2} \quad (\mathrm{mod}\ n)}

zu einem bestimmten Wert von $c\ (\mathrm{mod}\ n)$, der auch noch durch
8 teilbar angenommen werden kann. Wenn umgekehrt $x, y$
diesen Bedingungen genügen, so habe $\alpha$ und $\beta$ keinen gemeinsamen Teiler, wie aus

\[ \alpha\,\dfrac{Bx-y}{2} - \beta Cx = -n \]

hervorgeht. Denn danach müßte ein ungerader Teiler von $\alpha$
und $\beta$ in $n$ aufgehen, und müßte daher, da $A$ relativ prim zu $n$
angenommen war, in $x$ und folglich in $y$ aufgehen.

Es folgt aber noch, wenn wir von den beiden Ausnahmefällen
$D=-3$, $D=-4$ absehen, die uns hier überhaupt
nicht interessieren, weil für diese $j(\omega)$ rational ist, daß \gesperrt{verschiedene} Lösungen von (8) auch zu verschiedenen Werten $c$
führen. Denn nehmen wir an, daß ein und derselbe Wert von $c$
zu zwei verschiedenen Systemen $(\alpha, \beta, \gamma, \delta)$ führen könnte, so
würde aus (5) eine Relation der Form

\[ \omega = \dfrac{\gamma+\delta\omega}{\alpha+\beta\omega} \]

folgen, und diese Substitution ist, wenn sie nicht identisch ist,
nur für die beiden erwähnten Ausnahmefälle möglich.

\gesperrt{Demnach ist die Anzahl der $v$-Werte in (2), die nach
(3) gleich $u$ werden, so groß wie die Anzahl der eigentlichen Lösungen der Gleichung (8).}

Hat die Gleichung (8) nur \gesperrt{eine} solche Lösung, so ist der
entsprechende Wert von $M$ rational durch $u$ ausdrückbar. Hat
sie aber mehrere Lösungen, so sind die zugehörigen Werte von $M$
die Wurzeln einer rationalen Gleichung:

\weq{11}{\Phi(M,u) = 0,}

deren Grad in bezug auf $M$ ebenso groß ist, wie die Anzahl
dieser Lösungen (\S\,134).

Es genügt ferner $u$ der Klassengleichung

\[ H_{-D}(u) = 0. \]
